% TOP_PROBLEM: {"id": 2200010, "title": "Problems Around Polynomials — Conjecture 6", "category_id": 4, "category": {"id": 4, "name": "algebra", "display_name": "Algebra", "description": "Group theory, ring theory, field theory, and algebraic structures.", "slug": "algebra", "order_index": 4, "created_at": "2026-07-31T15:26:25.671Z"}, "status": "open"}
% FIRST_POSTED: null
% ENTRY_KIND: statement_only
% Imported from: batch09.tex; UnsolvedMath ID 2200010
\documentclass[11pt]{article}
\usepackage[margin=25mm]{geometry}
\usepackage{fontspec}
\setmainfont{Latin Modern Roman}
\usepackage{amsmath,amssymb,mathtools}
\usepackage{unicode-math}
\setmathfont{Latin Modern Math}
\usepackage{xurl,hyperref}
\hypersetup{colorlinks=true,urlcolor=blue}
\setcounter{secnumdepth}{0}
\setlength{\parindent}{0pt}
\setlength{\parskip}{5pt}
\setlength{\emergencystretch}{3em}
\newcommand{\sref}[2]{\hyperlink{source:#1:#2}{[S#2]}}
\newcommand{\cataloguescope}[1]{\paragraph{Recorded target.}#1}
\begin{document}
% UnsolvedMath ID 2200010; source catalogue code AMR-021-0010
\setcounter{section}{2200009}
\section{Problems Around Polynomials — Conjecture 6}\label{2200010}
{\small\sffamily UnsolvedMath ID 2200010 · Algebra}
\subsection{Definitions and mathematical statement}

\paragraph{Shared notation.}$\mathbb N=\{1,2,\ldots\}$, and $\mathbb Z,\mathbb Q,\mathbb R,\mathbb C$ denote the integers, rationals, reals and complex numbers. Any other conventions are specified locally.


For a nonzero real-rooted polynomial $p$ of degree at least two, its mesh is the minimum distance between consecutive real roots, listing roots with multiplicity. Thus a repeated root gives mesh zero. Nonzero constants and linear polynomials have mesh $+\infty$. As usual for linear real-rootedness preservers, the zero polynomial is an allowed zero output. Write $\mathcal H_m$ for zero together with real-rooted polynomials of degree at most $m$ and mesh at least one.\par For $d\geq1$, define the forward difference $\nabla p(x)=p(x+1)-p(x)$, with $\nabla^0p=p$. For $p,q\in\mathbb R[x]$ of degrees at most $d$, their finite-difference convolution is
\[(p\bullet_d q)(x)=\sum_{j=0}^{d}(\nabla^jp)(0)\,(\nabla^{d-j}q)(x).\]
\paragraph{Statement recorded in the supplied source.}
For every $d\geq1$ and all $p,q\in\mathcal H_d$, is $p\bullet_d q\in\mathcal H_d$? This is the convolution formulation of preservation of mesh at least one. \sref{2200010}{2} \sref{2200010}{3}
\subsection{Short English statement}
Does finite-difference convolution preserve real-rootedness and spacing at least one for polynomials of bounded degree?
\subsection{Sources}
\begin{description}
\item[\hypertarget{source:2200010:1}{S1}] ULAM.AI and the UnsolvedMath Contributors, UnsolvedMath: A Curated Collection of Open Mathematics Problems, record 2200010 (AMR-021-0010). CC BY 4.0. \url{https://huggingface.co/datasets/ulamai/UnsolvedMath}.
\item[\hypertarget{source:2200010:2}{S2}] Boris Shapiro, Problems Around Polynomials -- The Good, The Bad and The Ugly, Arnold Mathematical Journal 1 (2015), pp.91--104, Conjecture 6 and its complete preceding definitions. Original question and full discussion. \url{https://arxiv.org/abs/1503.05295}.
\item[\hypertarget{source:2200010:3}{S3}] Petter Brändén, Ilia Krasikov and Boris Shapiro, Elements of Pólya--Schur theory in finite difference setting (2013), Section 1 mesh definition and Section 4, Conjectures 2 and 3. \url{https://arxiv.org/pdf/1204.2963}.
\end{description}
\label{end:2200010}
\end{document}
